\documentclass{article}

\title{Partial Partitioning}
\date{}
\author{}

\begin{document}
\maketitle

\section{Mechanism}
Central to the halt-free implementation of partial partitioning is the idea of a logical clock for each partially-partitioned cluster, and when combined with the logical clocks of other clusters, a vector clock for the entire system.
These clocks are incremented by the coordinator of partitioned cluster and attached to every message.

\begin{enumerate}
  \item On the partitioned node, divide relations into partitionable and replicated. Any non-EDB relation that can't be partitioned (by the definition of the original paper) is labeled replicated.
  \item Any rule that has a replicated relation at the head has all relations in its body labeled replicated as well.
  \item All machines maintain a vector clock, one slot for each partitioned cluster.
  \item Upon receiving a message with a vector clock, non-partially partitioned machines update their own vector clock by taking the max of each slot.
  \item All machines attach their vector clock to all messages.
  \item Create a coordinator node (can be co-located with a partition).
    \begin{enumerate}
      \item Senders route all replicated inputs to the coordinator.
      \item When it receives an input tuple, it increments its cluster's slot in the vector clock and broadcasts to all partitions.
    \end{enumerate}
  \item Modify partitions:
    \begin{enumerate}
      \item It delays processing messages from senders until its cluster's slot in the vector clock is at least as high.
      \item It processes messages from the coordinator in incremental order, buffering until it receives the next message in the sequence.
    \end{enumerate}
\end{enumerate}

\section{Proof Sketch}
Original proof: Create a mapping function from the new partitioned states to the original, from a new $t'$ (on each of the partitions) to the original $t$, where $t' \ge t$ (so messages sent by the new nodes are delivered at a time they could've been delivered originally).

New proof: Shift times around based on the logical clock on each machine so that all states with a higher logical clock happen-after all states with a lower one.
Prove that this can be done consistently across all partitioned clusters.
Prove that there are no messages sent back in time as a result.



\end{document}
